\chapter{SDF Template and Registration}
\docsource{docs/design/05-sdf-pipeline.md}
\docsource{docs/design/06-registration.md}

\section{SDF Pipeline Boundary}

CloudCropper가 직접 소유하는 범위는 cropped point cloud에서 oriented-normal
template을 만들고, inspection PLY와 runtime NPZ를 내보내는 곳까지다. mesh
reconstruction과 gradient-SDF generation은 downstream Open3D/Python stack에 둔다.

\section{Normal Estimation}

normal 방향은 native \texttt{core}에서 PCA로 계산한다. 이미 계획된 nanoflann
KD-tree를 재사용하므로 별도 heavyweight dependency를 추가하지 않는다.

\begin{enumerate}
  \item denoise를 먼저 수행한다.
  \item \texttt{point\_spacing\_m} 기반 hybrid radius와 max neighbour count를
  선택한다.
  \item kNN 또는 hybrid neighbour set의 covariance를 만들고 가장 작은 eigenvalue의
  eigenvector를 normal direction으로 사용한다.
  \item sensor origin이 있으면 viewpoint flip으로 orientation을 정한다.
  \item sensor origin이 없고 전역 consistency가 필요하면 Open3D의 tangent-plane
  propagation에 위임한다.
\end{enumerate}

\section{Template Outputs}

\begin{longtable}{p{0.24\linewidth}p{0.66\linewidth}}
\toprule
\textbf{Output} & \textbf{용도} \\
\midrule
Inspection PLY & 사람이 열어 geometry와 normal을 확인한다. mesh/SDF 입력으로 바로
사용하지 않고, reconstruction stage의 입력 cloud로 둔다. \\
Runtime NPZ & downstream registration/SDF가 필요한 points, normals, bbox,
canonical transform, axis, spacing metadata를 schema로 전달한다. \\
\bottomrule
\end{longtable}

\section{Registration Backends}

registration은 cropped/saved cloud를 target으로, viewer에 열린 cloud를 source로
보며 target-from-source rigid transform을 추정한다.

\begin{lstlisting}[language=C++]
enum class RegAlgo {
  Icp,
  PlaneIcp,
  Gicp,
  VGicp,
  KissMatcher,
  KissGicp,
  GradientSdf
};

Result<RegResult> registerClouds(
  const PointCloud& source,
  const PointCloud& target,
  const RegOptions& opt = {});
\end{lstlisting}

공통 dispatcher는 backend 결과 이후 source-to-target nearest-neighbour RMSE와
inlier count를 다시 계산한다. 이 metric은 backend와 무관하게 비교 가능한
검증값이다.

\section{Viewer Panel}

Registration panel은 target path 또는 last export, algorithm combo, parameter
control, register button, result matrix, RMSE/inlier summary, apply/save/reset
actions를 갖는다. solve는 복사본 위에서 worker thread로 실행해 UI 상호작용을
막지 않는다.

